\section{Problem Formulation}
\label{sec:problem}

\subsection{State transition model}

Let $S_t=(M_t,\theta_t,\Graph_t,\sigma_t)$ denote the personalization state. Processing event $e_t$ applies a versioned transition
\begin{equation}
S_t=F_{v_t}(S_{t-1},e_t;\rho_t),
\label{eq:transition}
\end{equation}
where $v_t$ identifies code, model, tokenizer, hyperparameters, and numerical backend, and $\rho_t$ is an explicit randomness record. A memory transform $f_M$ may produce one or more derived records,
\begin{equation}
M_t=f_M(M_{t-1},e_t;\pi_{\mathsf{mem}}),
\label{eq:memory-update}
\end{equation}
while a trainer periodically computes
\begin{equation}
q_k=T(\theta_{k-1},W_k;v_k,\rho_k),\qquad
W_k\subseteq \{e_i\}\cup M_t.
\label{eq:training-update}
\end{equation}
The lineage graph records every membership and transformation relation needed to interpret these maps.

\subsection{Cross-sink laundering}

\begin{definition}[Cross-sink laundering path]
A laundering path for source $e$ is a directed path
\begin{equation}
e\leadsto m\leadsto w\leadsto q\leadsto \theta
\label{eq:laundering-path}
\end{equation}
where $m$ is a memory or derived record, $w$ is a training example/window, $q$ is a candidate adapter, and $\theta$ is a promoted checkpoint. The path is successful when deleting $e$ and $m$ leaves trigger-conditioned behavior in $\theta$ above a declared success threshold.
\end{definition}

The empirical laundering score separates persistence from initial compromise. For trigger distribution $\mathcal{T}$, malicious target predicate $g$, active state $S$, and post-deletion state $S^{-e}$,
\begin{align}
\mathrm{ASR}(S)&=\Pr_{z\sim\mathcal{T}}[g(f_S(z))=1],\\
\mathrm{LP}(e)&=\mathrm{ASR}(S^{-e})-\mathrm{ASR}(S_{\mathrm{clean}}).
\label{eq:laundering-metric}
\end{align}
$\mathrm{LP}$ is not inferred from a single example; it is estimated across triggers, seeds, and workloads with confidence intervals after actual experiments.

\subsection{Utility and systems objectives}

Let $D_{\mathrm{clean}}$ be a held-out benign personalization set and $U(S;D_{\mathrm{clean}})$ its task-appropriate utility. False quarantine rate is
\begin{equation}
\mathrm{FQR}=\frac{|\{e\in D_{\mathrm{benign}}:\pi(e)=\Quarantine\text{ or }\Deny\}|}{|D_{\mathrm{benign}}|}.
\end{equation}
The defense aims to reduce attack risk subject to explicit utility and resource budgets:
\begin{equation}
\min_{\pi,A,I} \ \Risk_{\mathrm{attack}}+\lambda_U\Delta U+\lambda_Q\mathrm{FQR}
\quad\text{s.t.}\quad C_S\le B_S,\ C_E\le B_E,\ C_T\le B_T,
\label{eq:objective}
\end{equation}
where $A$ is the audit gate, $I$ is the checkpoint interval, and the $\lambda$ values express deployment preferences rather than universal constants.

\subsection{Design requirements}

The problem produces six requirements. R1: preserve origin and derivation across transformations. R2: mediate each sink separately. R3: make unpromoted adapters non-serving. R4: connect deletion/revocation to every consuming checkpoint. R5: restore from a known unaffected state rather than trusting approximate parameter surgery. R6: expose storage, latency, energy, and benign-content costs on phones.

\subsection{Non-objectives and impossibility boundary}

No content-only policy can guarantee that an authenticated user-authored example reflects an unmanipulated preference. If the adversary persuades the user to type the exact poisoned example, the provenance tuple legitimately changes. PoisonGuard can require confirmation for high-impact learning or rate-limit repeated patterns, but cannot infer intent perfectly. Similarly, if the TCB is compromised, the attacker can forge or suppress edges. These are model boundaries, not implementation footnotes.

\begin{figure}[t]
  \centering
  \resizebox{0.95\columnwidth}{!}{\begin{tikzpicture}[
  hnode/.style={honest, font=\small, text width=31mm, minimum height=10mm, inner sep=3pt},
  dnode/.style={defense, font=\small, text width=31mm, minimum height=10mm, inner sep=3pt},
  inode/.style={infra, font=\small, text width=31mm, minimum height=10mm, inner sep=3pt},
  anode/.style={attack, font=\small, text width=31mm, minimum height=10mm, inner sep=3pt}
]
  % Row 1: Ingress & Policy Selection (y = 0)
  \node[hnode] (s1) at (0, 0) {\textbf{1. Capture OS Origin}\\[0.3mm]\scriptsize Provenance envelope binding};
  \node[dnode] (s2) at (5.2, 0) {\textbf{2. Select per Sink}\\[0.3mm]\scriptsize Policy rules ($\pi_{\text{mem}}, \pi_{\text{train}}$)};

  % Row 2: Sink Ingestion & Isolated Training (y = -2.0) -> Vertical gap = 10mm!
  \node[hnode] (s3) at (0, -2.0) {\textbf{3. Buffer in Memory}\\[0.3mm]\scriptsize Memory records \& windows};
  \node[inode] (s4) at (5.2, -2.0) {\textbf{4. Quarantine Train}\\[0.3mm]\scriptsize Isolated LoRA adapter};

  % Row 3: Audit Promotion & Rollback Remediation (y = -4.0) -> Vertical gap = 10mm!
  \node[dnode] (s5) at (0, -4.0) {\textbf{5. Audit \& Promote}\\[0.3mm]\scriptsize Lineage closure gate};
  \node[dnode] (s6) at (5.2, -4.0) {\textbf{6. Lineage Rollback}\\[0.3mm]\scriptsize Causal rollback \& replay};

  % Long, prominent 21mm horizontal arrows & 10mm vertical arrows (Matching Fig 2 style)
  \draw[dataflow] (s1.east) -- (s2.west);
  \draw[dataflow] (s1.south) -- (s3.north);
  \draw[dataflow] (s2.south) -- (s4.north);
  \draw[dataflow] (s3.east) -- (s4.west);
  \draw[controlflow] (s3.south) -- (s5.north);
  \draw[controlflow] (s4.south) -- (s6.north);

  % 21mm long revocation trigger arrow with matching Fig 2 compact orange style
  \draw[attackflow] (s6.west) -- node[midway,above=1mm,font=\scriptsize\bfseries,text=pgorange,fill=white,inner sep=1.2pt,rounded corners=1pt]{on revocation} (s5.east);
\end{tikzpicture}
}
  \caption{PoisonGuard method pipeline. Numbered stages correspond to capture, sink-specific selection, derived-state construction, quarantined training, audit/promotion, and revocation-driven rollback.}
  \label{fig:pipeline}
\end{figure}
